% !TeX program = xelatex
% ==========================================================================
% Звіт з лабораторної роботи — ДВВ «Основи прикладної фізики» (1.26-27), УКУ.
% Викладач: Володимир Коломієць.
% Це ЄДИНИЙ файл, який вам потрібно редагувати. Уся верстка й брендинг
% живуть у ucaphyslab.cls — не змінюйте шрифти, поля, кеглі й кольори.
%
% Компіляція: XeLaTeX + BibTeX (Overleaf робить це автоматично).
% Локально:   latexmk (див. README.md у корені репозиторію).
% ==========================================================================
\documentclass[ukrainian]{ucaphyslab}
% Опції класу (можна комбінувати через кому):
%   ukrainian | english  — мова документа й логотипу (за замовчуванням ukrainian)
%   slogan                — додає гасло УКУ в колонтитул титулки (за замовчуванням вимкнено)

% --------------------------------------------------------------------------
% Метадані титульної сторінки
% --------------------------------------------------------------------------
\labtitle{Маятник Обербека}
\labsubtitle{Лабораторна робота № (вкажіть номер)}

\student{Прізвище Ім'я}
% \student{Прізвище Ім'я} % другий співвиконавець (опційно, для групової роботи)

\teacher{Володимир Коломієць} % якщо рядок не потрібен — видаліть текст і залиште \teacher{}

\labdate{РРРР-ММ-ДД}

\begin{document}
\maketitlepage

\section{Мета роботи}
Сформулюйте, що саме досліджується і яку фізичну величину чи закономірність
потрібно визначити чи перевірити.

\section{Прилади та обладнання}
\begin{itemize}
  \item Перелічіть використані прилади, їх клас точності та ціну поділки.
\end{itemize}

\section{Теоретичні відомості}
% Основні формули, з яких обчислюються шукані величини, з короткими
% поясненнями позначень.
Наведіть теоретичні відомості та розрахункові формули, наприклад:
\begin{equation}
  T = 2\pi\sqrt{\frac{l}{g}}
\end{equation}

\section{Хід роботи та результати вимірювань}
% Опишіть послідовність дій. Результати прямих вимірювань зводьте в таблицю.
Опишіть послідовність виконання роботи.

\begin{table}[h]
  \centering
  \begin{tabular}{@{}cccc@{}}
    \toprule
    № & Величина 1 & Величина 2 & Величина 3 \\
    \midrule
    1 & & & \\
    2 & & & \\
    3 & & & \\
    \bottomrule
  \end{tabular}
  \caption{Результати вимірювань}
\end{table}

\section{Обробка результатів}
% Розрахунки шуканих величин та оцінка похибок вимірювань.
Наведіть розрахунки та підсумковий результат разом із похибкою, наприклад:
\begin{equation}
  g = (\num{9.8} \pm \num{0.1})\ \si{\meter\per\second\squared}
\end{equation}

% Приклад врізки з ключовим результатом (опційно, можна видалити):
% \begin{keyfacts}
% \textbf{Результат:} $g = (9.8 \pm 0.1)\ \text{м/с}^2$ \textperiodcentered\ відносна похибка 1{,}0\%
% \end{keyfacts}

\section{Висновки}
Сформулюйте висновок: чи узгоджується отриманий результат із теоретичним
(табличним) значенням у межах похибки, та які фактори могли вплинути на
точність вимірювань.

% Розкоментуйте, якщо цитуєте джерела через \cite{} (записи — у references.bib):
% \ucuprintbibliography

\end{document}
